The starting point is the contrast between the two stacks at low frequency. At 1~kHz, both devices show the butterfly-shaped permittivity characteristic of ferroelectric \hzo{}, and both reach a similar peak relative permittivity of about 62. Yet the zero-bias state separation is different: Stack~A gives $\CMW_\xi \approx 5.2$, whereas Stack~B gives only $\approx 2.3$ [Fig.~\ref{fig:cmw}]. This immediately suggests that the zero-bias read margin is not set only by the ferroelectric dielectric amplitude. As the measurement frequency increases, the branch separation collapses while a substantial background permittivity remains. The disappearing quantity is therefore the state discrimination itself, not the dielectric response as a whole.

The frequency dependence of that state discrimination is captured by a single-Debye form,
\begin{equation}
\CMW(f)=\frac{\CMW_0}{1+(f/\fc)^2},
\label{eq:debye_short}
\end{equation}
which directly isolates the relaxing part of the zero-bias window. Fits to Fig.~\ref{fig:debye} give $\CMW_0=5.38$ and $\fc=179$~kHz for Stack~A, corresponding to $\tau=0.90\,\mu$s, while Stack~B relaxes much more slowly with $\CMW_0=3.6$ and $\fc=1.4$~kHz ($\tau=114\,\mu$s). The full TMAT dataset shows that the differential loss and Cole--Cole representations return the same characteristic corner and do not require a distributed relaxation spectrum. That behavior is consistent with one dominant relaxation process rather than with a broad intrinsic ferroelectric response~\cite{jonscher1983,colecole1941}.

Temperature provides an independent discriminator. At 10~K, Stack~A still switches: $2P_r$ decreases from 33.1 to 20.3\uCcm{} and $E_c$ increases from 0.93 to 1.17\MVcm{}, but the device remains ferroelectric [Fig.~\ref{fig:cryo}]. The capacitive readout changes more severely. The zero-bias CMW falls from 5.2 to about 2.2 at 1~kHz and then drops to the noise floor above 1.1~kHz, implying $\fc \lesssim 1$~kHz at 10~K. Cooling therefore suppresses the capacitive-read bandwidth by more than two decades even though polarization reversal remains observable. That decoupling is difficult to reconcile with a read mechanism controlled only by the ferroelectric switching timescale.

The strongest argument comes from comparing the two stacks directly. Their low-frequency dielectric ceilings are nearly identical, so the ferroelectric layer presents essentially the same dielectric amplitude in both devices. Nevertheless, the relaxation times differ by $126\times$ [Table~\ref{tab:summary}], and the stack containing the deliberate WO$_x$ interface is the slower one. A compact Maxwell--Wagner description,
\begin{equation}
\tau \approx \frac{\varepsilon_i}{\sigma_i},
\label{eq:mw_short}
\end{equation}
assigns that timescale to the interfacial layer rather than to the ferroelectric bulk~\cite{alcala2023,kremer2003}. The extracted interfacial areal capacitances differ only modestly (2.95 versus 2.34~$\mu$F/cm$^2$), whereas the corresponding areal resistances differ by about $157\times$. In other words, the two stacks share nearly the same dielectric ceiling but not the same interface clock.

Taken together, the full plot set supports a simple reading rule for the reduced manuscript. The zero-bias CMW exists because polarization creates two stored states, but the rate at which those states remain capacitively distinguishable is controlled by interfacial charge relaxation. The metric reported at a single frequency is therefore not portable by itself. For conference presentation, the critical outputs are the low-frequency window, the fitted corner frequency, and the interface-sensitive comparison that shows why the same ferroelectric can yield very different NDRO read bandwidths.
